% Transcription — fonds Alexandre Grothendieck, shelfmark 69, batch 5 (pages 81-100).
% Established from the facsimile archives/batches/69/batch-05.pdf (cut from a
% copy of 69.pdf fetched through the project's relay,
% https://grothendieck-relay.onrender.com/source/69.pdf), per the skill
% transcribe-grothendieck. The folder is 119 pages in six batches (1-20,
% 21-40, 41-60, 61-80, 81-100, 101-119), transcribed in parallel; this is
% batch 5. The raw scan's cover sheet is removed from the batch PDFs, so PDF
% page k is page 80+k, which is also the Montpellier footer number.
%
% Pass: Opus 5.5 (claude-opus-5-5), 2026-09-27 — first pass, unchecked against the pages by a human.
%
% Numbering. \page{} follows the PDF count (80+k), as folder 121 does. The
% archivists' pencil numbers bottom left agree with it on PDF pages 1-13
% (« 81 » ... « 93 »), then run two ahead: PDF pages 14-20 are pencilled
% « 96 » ... « 102 » (no « 94 », « 95 » on these leaves). Pages 95-100
% carry a \note giving their pencil number. He paginates nothing here.
%
% Pages skipped: 88, 93, 94 and 99, blank cream separator sheets (only
% show-through). Page summarised: 83, a typed announcement (USTL,
% D.E.U.G. A première année, option mathématique, 1976, « Polyèdres
% réguliers », a course in two parts by A. Grothendieck and C. Contou-Carrère)
% on the recto of the leaf whose verso is page 82; nothing of his hand on it,
% one French \note, following transcripts/71/batch-01.fr.tex. This is
% presumably the typed item the contact sheet flagged; it is not a letter.
% Its « 1976 » is consistent with the archivists' « [à partir de 1976] ».
%
% Pages transcribed: 81, 82, 84-87, 89-92, 95-98, 100, all in his hand,
% blue ink, on plain paper (81-83) and on punched computer-listing paper.
% Much is fast drafting; the formulas carry better than the connective
% prose. Pages 90 and 91 are cancelled (a cross; wavy strokes) but legible,
% and are transcribed with the cancellation said once. The drawings (a
% pentagonal fragment on 81, stars and « mobiles » of labelled vertices on
% 84, 86, 87, strips of triangles on 85) are described in \drawing{}, not
% redrawn: none is simple enough to redraw exactly.
%
% What the batch is about. Apparently not graphs as such but the configuration
% of the 27 lines on a cubic surface and the root system E6 (72 roots,
% |W| = 27·10·8·4·6 = 51840), though he never names them here.
%   81-87  Combinatorics of the 27 : vertices ξ_i, ξ'_i, ξ_ij (i, j in 1..6),
%          « liés » / « non liés » pairs; counting 27·16, 27·10, 72 roots
%          (« de carré 2 »), 270, 216, 243 elements (82); indices of
%          stabilisers in the Weyl group, [W] : n, for pinned / pointed /
%          oriented prisms, A_5, A_4, A_3, « carrousels », hexagons,
%          pentagons, triangles (84, 86); the pairing δ·δ' ≡ card(δ∩δ') +
%          card(ε(δ)∩ε(δ')), and x, y liés iff card u = 3 or 4 for
%          u = x - y in F_2^A (86); « bicarrousels », 1080 = 72·15 (87).
%   89-92  The 72 roots: relative positions of two roots (equal, opposite,
%          orthogonal, 60°, 120°; 1, 1, 30, 20, 20), 120 special planes (89);
%          reflections s_r(x) = x + (r,x) r, and the vectors
%          -ξ~_i = -(3ξ_i - λ), -ξ~'_i, -η~ expressed in simple roots
%          r_1 ... r_6 (r_0), with a table (90-92); solving for
%          x = aη + bξ ± ξ_i (± ξ_j) and the symmetry δ -> 2λ_0 - δ,
%          r_i -> -r_{6-i} (91).
%   95-98  For a root system R ⊂ E with Weyl group W: the sets 𝔖 (x with
%          (x,r) in {-1,0,1} for all r) and 𝔖_0 (wx - x in R whenever
%          wx ≠ x), card 𝔖 ≤ 3^r - 1; the cases A_r (two W-orbits for
%          r ≥ 2, exchanged by the symmetry), B_r, C_r (𝔖_0 empty), D_r
%          (three orbits, 𝔖_0 empty); then 𝔖'_0, 𝔖' with a ray condition,
%          and the conclusion that 𝔖'_0 = 𝔖_0 ∪ R exactly when all roots
%          have the same length (A_r, D_r, E_6-E_8), with the exception
%          W.(-½,-½,½,½) for A_3.
%   100    Unrelated group theory in his hand: three normal subgroups
%          N_1, N_2, N_3 pairwise meeting in {e} and pairwise generating G,
%          N_1 × N_2 ≃ G, G ≃ N_2 ⋊ N_1 ≃ N_3 ⋊ N_1. Transcribed.
%
% Notation, fixed once. His looped capital (a curly S) for the sets of
% pages 95-98 is set \mathfrak{S}; the tilde vectors of 90-92 are
% \tilde\xi, \tilde\eta; his « ép. » is kept as written (« épinglé »).
% His circled numerals are set (n) with a \note saying they are circled.
%
% The dating is the folder's own; the group « Géométrie et topologie
% combinatoire » (69-89) holds it.
\documentclass[11pt,a4paper]{article}
\input{../preamble/grothendieck}

\folder{69}
\batch{5}
\pages{81}{100}
\dating{[à partir de 1976]}
\watermark{Édition de démonstration}
\foldertitle{Graphes cubiques : notes manuscrites (s.d.), lettre (s.d.).}

\begin{document}

\page{81}
\drawing{Un fragment de graphe cubique à faces pentagonales (six
pentagones autour d'un pentagone central, à la manière d'un dodécaèdre
déplié). Les arêtes portent les numéros 1 à 5 ; les trois arêtes
numérotées 3 qui forment un couplage sont repassées en gras. Les sommets
sont marqués : au pentagone central, $\xi_{12}$, $\xi_{45}$, $\xi_{23}$,
$\xi_{51}$, $\xi_{34}$ ; à l'extérieur, $\xi_{35}$, $\xi_{4,12}$
\uncertain{}, $\xi_{13}$, $\xi_{25}$, $\xi_{41}$, $\xi_{5,3}$,
$\xi_{2,4}$, $\xi_{3,1}$, $\xi_{52}$, $\xi_{14}$. En haut à gauche, un
petit pentagone raturé, une de ses arêtes en gras.}
\note{En haut à droite, isolé, le nombre « 156 » (lecture douteuse :
peut-être « 146 »).}

\page{82}
\[
  \delta'_1-\delta_1=\delta'_2-\delta_2 \qquad
  \delta_1\neq\delta_2,\ \delta_1\neq\delta'_1
\]
\[
  \underbrace{\delta'_1+\delta_2}=\underbrace{\delta'_2+\delta_1}
\]
$\delta'_1$ lié à $\delta_2$, $\delta'_2$ lié à $\delta_1$ ;
$\delta'_1$ lié ni à $\delta'_2$ ni à $\delta_1$.
\[
  \xi_{ij}=\eta-\xi_i-\xi_j \qquad
  \sum\xi_{ij}=15\eta-5\xi=5\lambda
\]

\emph{Différences de sommets}\note{Le mot est souligné. Les quatre
alinéas qui suivent sont réunis par un trait vertical et ouverts chacun
par un signe cerclé, qu'on lit « 0 », « 6 », « 1 » \uncertain{} et
« b » \uncertain{}.}

(0) sommets non liés se groupent par paquets de 6 parmi les $27\cdot16$
possibilités initiales, d'où $72$ « racines » (de carré $2$) ;

(6) sommets liés donnent $27\cdot10=270$ autres éléments
distincts\add{ correspondant aux arêtes}, \ill{} \uncertain{(de carré
$-4$)}, \ill{} en tout $342$\note{Lecture douteuse : on peut lire aussi
« 392 » ; $72+270=342$.} éléments (de carré \ill{}) ;

(1) sommets non liés (distincts) se groupent par paquets de $5$ parmi
les $27\cdot10/2=135$ possib.\ initiales, \ill{} d'où $27$ \ill{} (de
carré $2$) correspondant aux $27$ \struck{triangles} sommets \ill{}
\ill{} les \struck{\ill{}} triangles adjacents, et les autres \ill{}
\[
  \delta\longmapsto\lambda-\delta
\]

(b) \struck{\uncertain{pointés}} sommets non liés (\add{correspondant aux
\uncertain{bicycles} pointés}) donnent $27\cdot16/2=216$ autres
sommets distincts, formant avec les précédents $27+216=243$ éléments.
\[
  \xi_i\longmapsto\underbrace{\xi'_j,\ \xi_{ij}}\qquad
  \xi'_j-\xi_{ij}=2\eta-\xi+\xi_h
\]
\[
  \struck{=\eta-\xi_i-\xi_j}\qquad \struck{\ill{}}\ -\xi_i
\]
\[
  \xi'_i\longmapsto\underbrace{\xi_j,\ \xi_{ij}}\qquad
  \xi_h+\xi_{ij}=\eta-\xi_j
\]
\note{La ligne biffée du milieu est incertaine : les deux dernières
égalités sont écrites en marge d'un trait vertical et les indices $h$
se lisent mal.} \ill{} \ill{} ou $3$ sommets

\page{83}
\note{Page dactylographiée, recto de la feuille dont la page 82 est le
verso ; rien de sa main n'y figure. C'est l'annonce d'un enseignement de
l'Université des Sciences et Techniques du Languedoc, Mathématiques,
D.E.U.G.\ A première année, option mathématique, 1976, intitulé
« Polyèdres réguliers » : un « cours » pour étudiants « disposés à
trouver de l'amusement en faisant des choses mathématiques très
élémentaires et "visuelles" », sans travaux dirigés ni exercices notés,
en deux parties développées l'une par A.~Grothendieck, l'autre par
C.~Contou-Carrère. Programme : I) polyèdres — polygones plans, division
du plan en deux parties (« th.\ de Jordan »), cas convexe, polyèdres
convexes et leurs facettes, combinatoire d'un polyèdre convexe, dualité
des polyèdres convexes ayant un point intérieur donné ; II) les cinq
polyèdres réguliers (« solides platoniciens ») — descriptions métrique,
combinatoire, éventuellement « projective » en termes de configurations
de points, droites et plans, étude de leurs groupes de symétrie et liens
avec les groupes symétriques $\mathfrak{S}_3$, $\mathfrak{S}_4$,
$\mathfrak{S}_5$ et leurs sous-groupes. Les lettres des groupes
symétriques semblent tracées à la main dans les blancs de la frappe.}

\page{84}
\note{Feuille de listing tenue en largeur. En haut à gauche, une
chaîne d'équivalences et d'inclusions, qu'on lit de bas en haut.}
\[
\begin{array}{c}
  A_3\ \text{ép.}\\
  \uparrow{\scriptstyle 4}\\
  A_4\ \text{ép.}\sim A_5\ \text{ép.}\sim\text{hex.\ ép.}\sim
  \text{prisme épinglé}\sim\text{carrous.\ ép.}\\
  \uparrow{\scriptstyle 3}\\
  \text{carrousel ép.}\sim\text{pent.\ ép.}\\
  \uparrow{\scriptstyle 2}\\
  (\text{\uncertain{repère}})
\end{array}
\]
\note{« ép. » abrège « épinglé » ; « carrousel » est lu d'après
« bicarrousel », écrit en clair à la page 87.}

\drawing{Au milieu de la feuille, un grand dessin : un triangle dont
les trois sommets sont marqués $\xi_{56}=\omega$ (en haut), $\xi_{12}=$
\ill{} (en bas à gauche) et $\xi_{34}=$ \ill{} (en bas à droite) ; de
chacun part un éventail d'arêtes vers des sommets marqués d'indices
doubles ou simples, primés ou non, certains cerclés, noircis ou encadrés.
À gauche, les extrémités sont groupées par des accolades $a_1$, $a_2$,
$a_3$, $a_4$ et portent notamment $\xi_1=\alpha_4$ \uncertain{},
$\xi'_2=\alpha'_1$, $\xi_{35}$, $\xi_{46}=\alpha'_3$ \uncertain{},
$\xi_{36}$, $\xi_{45}=\alpha_4$ ; en haut, les accolades $c_1$, \ldots,
$c_4$, avec $\xi_{23}$, $\xi_{14}$, $\xi_{13}$, $\xi'_3$, $\xi_6$,
$\xi'_6$, $\beta_1=\xi_5$ ; à droite et en bas, les accolades $b_1$,
\ldots, $b_4$, avec $\beta_1$, $\beta_2$, $\beta_3$, $\beta_4$,
$\xi_{51}$, $\xi_{52}$, $\xi_{26}$, $\xi_3$, $\xi_4$, $\xi'_3$,
$\xi'_4$. Beaucoup d'indices, écrits serré et surchargés, se lisent
mal.}

Sous le dessin, à gauche :
\[
  a_1+a_4=\alpha_1,\qquad a_2+a_4=\alpha_2,\qquad a_3+a_4=\alpha_3
\]

\drawing{À droite, en haut, un prisme triangulaire de sommets $a$, $b$,
$c$, $a'$, $b'$, $c'$, quatre arêtes repassées en gras ; au-dessous, des
lignes brisées de cinq, de trois et de deux points. Plus bas, deux
croquis : un triangle posé sur un rectangle ouvert par le bas, avec une
hauteur tracée ; un quadrilatère surmonté d'un triangle, ses sommets
inférieurs reliés en croix.}

À droite, en colonne :
\[
\begin{array}{ll}
  \text{\emph{prismes épinglés}}\ \ 27\cdot10\cdot8\cdot4\cdot1=[W]:6 &\\
  \text{prismes} & [W]:72\\[4pt]
  \text{A}_5\ \text{épinglés}\ \ 27\cdot10\cdot8\cdot4\cdot1=[W]:6 &\\
  \text{A}_5 & [W]:12\\[4pt]
  \text{A}_4\ \text{épinglé}\ \ 27\cdot10\cdot8^{\,4}=[W]:\uncertain{6} &\\
  \text{A}_4 & [W]:12\\[4pt]
  \text{A}_3\ \text{épinglés}\ \ 27\cdot10\cdot8=[W]:24 &\\
  \quad(\text{autres que les triangles épinglés}) &\\
  \text{A}_3\ \text{ép.} & [W]:48\\[4pt]
  \text{carrousels}\ \text{\add{simples}}\ \text{épinglés}\ \
  27\cdot10\cdot8\cdot4\cdot3=[W]:2 &\\
  \text{carrousels}\ \text{\add{simples}} & [W]:12
\end{array}
\]
\note{Dans « A$_4$ épinglé », un « 4 » est écrit au-dessus du $8$
(facteur ajouté ou exposant), le $6$ est surchargé, et le $12$ de la
ligne suivante est repassé sur un autre chiffre ;
dans « A$_3$ épinglés », l'expression entre crochets est surchargée et
se lit mal : $[W]$ est la lecture probable. Les titres « prismes
épinglés » et « carrousels \ldots\ épinglés » sont soulignés et reliés
par un trait à leur ligne.}

\page{85}
\drawing{Feuille de listing, presque vide, dessins à l'encre bleue
tracés la feuille tournée. En haut, deux bandes de triangles accolés
(des « anneaux » de triangles, comme la ceinture d'un antiprisme),
leurs sommets numérotés $1,2,3,4,5$ d'un côté et $1',2',3',4',5'$ de
l'autre, certaines arêtes repassées en gras ; à côté, un carré avec ses
deux diagonales. Au-dessous, une troisième bande analogue, sommets
$1'',2',3',4',5'$ \uncertain{}. En bas, un éventail : un triangle et
quatre arêtes issues d'un même sommet, avec des marques d'indices
illisibles.}

\page{86}
\[
  \xi_1\quad\xi_{14}\quad\xi_4\quad\xi'_2 \qquad\qquad
  \xi_1\ \big|\ \xi'_2\quad\xi'_3\quad\xi_{31}\quad\xi_{12}
\]
\drawing{En haut à droite, un croquis de polyèdre en perspective
(prisme ou antiprisme), quelques arêtes repassées.}
\drawing{À gauche, des segments étiquetés $\delta$, $\delta'$,
$\delta''$ (tracés en zigzag), puis $\delta'''$, $\delta^{IV}$.}
\[
  \delta=\xi_1,\ \delta'=\xi_2,\ \delta''=\xi'_2 ;\qquad
  \delta'''=\xi_{14},\ \delta^{IV}=\xi_{25}
\]
\[
  \delta\delta'=0\quad\delta\delta''=1\quad
  \delta\delta'''=1\quad\delta\delta^{IV}=0
\]
\note{L'indice de $\delta'$ est surchargé.}
\[
  \delta\delta'\ \struck{\bmod 2}\equiv
  \operatorname{card}(\delta\cap\delta')
  +\operatorname{card}\bigl(\varepsilon(\delta)\cap\varepsilon(\delta')\bigr)
\]
\drawing{Une rangée de croquis : un segment pointé en son milieu ; un
éventail de trois traits ; deux barres qui encadrent une croix
($|\times|$) ; trois barres parallèles ; deux barres et un triangle
allongé ; une barre, une diagonale, une barre. Au-dessous, à gauche, un
triangle et un éventail dont les sommets sont marqués $x$, $1$, $2$,
$3$, $3'$, $4$, et un hexagone de sommets $\xi_3$, $\xi_{23}$, $\xi_2$,
$\xi_{12}$, $\xi_1$, $\xi_{31}$, repassé, avec des cordes ; à côté, un
pentagone.}
\[
\begin{array}{ll}
  \text{hex.\ pointés \add{orientés}}\ \
  27\cdot10\cdot8\cdot4\cdot1\cdot1=2^6\cdot3^3\cdot5=[W]:6 &\\
  \text{hexagones} & [W]:72\\[4pt]
  \text{pentagones pointés \add{orientés}}\ \
  27\cdot10\cdot8\cdot4\cdot3=2^6\,3^4\,5=[W]:2 &\\
  \text{pentagones} & [W]:20\\[4pt]
  \text{carr.\ pointés}\ \ 27\cdot10\cdot8\cdot4=[W]:6 &\\
  \text{carr.} & [W]:48
\end{array}
\]
\note{« orientés » est ajouté au-dessus de la ligne aux deux premières
entrées. Le $6$ de la première ligne et le $72$ de la deuxième sont
repassés sur d'autres chiffres ; sous « $27\cdot10\cdot8\cdot4\cdot1
\cdot1$ », de petites accolades groupent les facteurs deux à deux.}
\drawing{À gauche, une étoile marquée $A$, et au-dessus
$\mathbb{F}_2^{\,i}$ \uncertain{} et un trait vertical.}
\[
  \mathbb{F}_2^{A}\quad P\qquad x,y\qquad
  x-y=u\in\mathbb{F}_2^{A}\simeq\mathfrak{P}(A)
\]
formé des $\alpha\in A$ sur lesquels \ill{} \ill{}
\[
  x,y\ \text{liés}\iff
  \begin{cases}
    x\ \text{et}\ y\ \text{de même parité}\ (\text{i.e.\ card}\ u\
      \text{pair}) :\ x-y=\mathbb{1}^{A}\\
    \quad(\text{i.e.}\ E(x)\cap\tilde A\ \text{et}\ E(y)\cap\tilde A\
      \text{sont complém.})\\
    x\ \text{et}\ y\ \text{de parité contr.}\ (\text{i.e.\ card}\ u\
      \text{impair}) :\ \struck{\ill{}}\ \operatorname{card}u=3
  \end{cases}
\]
\note{En face, reliés à la première ligne par une accolade,
« $E(x)\cap\tilde A$ » et « $E(y)\cap\tilde A$ », suivis de « sont
distincts » biffé ; devant « $x-y$ », un mot biffé \ill{}.}
\[
  x,y\ \text{non liés}\iff
  \begin{cases}
    x\ \text{et}\ y\ \text{de même parité i.e.\ card}\ u\ \text{pair} :\
      x-y\neq\mathbb{1}\\
    \quad(\text{i.e.}\ \struck{\ill{}}\ E(x)\cap\tilde A\cap
      E(y)\cap\tilde A\neq\emptyset)\\
    x\ \text{et}\ y\ \text{de parité contr.} :\ \operatorname{card}u=1
  \end{cases}
\]
\note{Sous la dernière ligne, un complément illisible commençant par
« i.e.\ ».}
\drawing{À gauche, au milieu, la mention « Conditions sur $\xi\in\tilde
A$ » et « $\xi$ liés à \ill{} », avec $\lambda_{12}$, $\lambda_{23}$,
$\lambda_4$ ; au-dessous, une étoile dont le sommet central est relié à
des sommets marqués $\lambda_{56}$, $\xi_1$, $\xi'_1$, $\xi_2$,
$\xi'_2$, $\lambda_{34}$, $\lambda_{12}$, $\lambda_{35}$, $\lambda_{46}$,
$\lambda$ \uncertain{}. En bas, des segments parallèles joignant des
points à des croix.}
\[
  x,y\ \text{liés}\iff\operatorname{card}u=3\ \text{ou}\ 4 \qquad
  x,y\ \text{non liés}\iff\operatorname{card}u=0,1\ \text{ou}\ 2
\]
\note{Sous le « $0$ », un renvoi : « $x=y$ ».}
\[
\begin{array}{l}
  \text{triangles pointés orientés}\ \ 27\cdot10=270=[W]:
    \uncertain{2^3\cdot3^4}\\
  \text{triangles}:\ \dfrac{27\cdot10}{6}=45=[W]:\uncertain{2^7\cdot3^2}
\end{array}
\]
\note{Les exposants sont repassés et se lisent mal ; ceux de la
première ligne sont rejetés dans la marge.}

\page{87}
\drawing{En haut, quatre petits rectangles ouverts d'un côté, points et
croix aux coins ; à gauche, un « i » et un croquis de triangles et de
quadrilatère, un éventail en bas, quelques sommets marqués d'une croix.
Au milieu, un « mobile » : un triangle de sommet $\xi_{56}$, de base
$\xi_{12}$--$\xi_{34}$, d'où pendent des arêtes vers $\xi_1$, $\xi_2$,
$\xi_6$, $\xi_5$, $\xi_3$, $\xi_4$, avec au centre, cerclé,
$\xi_{24}$ \uncertain{}. Plus bas, $\xi_{46}$ et $\xi_{62}$
\uncertain{} (surchargé), et trois éventails de trois branches.}
\[
\begin{array}{l}
  \text{bicarrousel épinglé}\\
  \text{\ill{} du bicarrousel}\\
  \quad\text{d'indice}\ 48=2^4\cdot3\\
  \text{nb des bicarrousels}\ \ 2^3\,3^3\,5=6^3\cdot5=1080\\
  \quad=\underbrace{72}_{\text{\uncertain{racines}}}\cdot
    \underbrace{15}_{\text{nb de \ill{} d'une \uncertain{racine}, 6}}
\end{array}
\]
\note{Le $15$ est repassé en gras.}
$\xi$ \ill{} \ill{} des \uncertain{paires} \ill{} $i>\ell$ \uncertain{}
\drawing{Un petit croquis : des sommets $\xi_1$, $\xi_4$, $\xi_{62}$ et
d'autres marqués d'indices illisibles, avec deux arêtes en croix, en
partie biffé.}
\[
  \frac{8!}{2!\,\struck{\ill{}}\,2!}\cdot
  \binom{6}{2}\binom{4}{2}\,\frac{1}{3!}
  =\frac{15\cdot6}{6}=15
\]

\page{89}
\emph{Positions relatives de deux racines} (ou de deux bases)

Soient $r,r'$ deux racines, correspondant à des bases \struck{\ill{}}
$b,b'$
\[
\begin{array}{lll}
  r'=r & \iff r\cdot r'=-2 & \iff b=b'\\
  r'=-r & \iff r\cdot r'=2 & \iff b'\ \text{associé à}\ b\\
  r'\ \text{orth.\ à}\ r & \iff rr'=0 & \iff \operatorname{card}(b\cap b')=1\\
  \widehat{r,r'}=60^\circ & \iff rr'=-1 & \iff
    \operatorname{card}(b\cap b')=3\ \ \struck{\text{\ill{}}}\\
  \widehat{r.r'}=120^\circ & \iff rr'=1 & \iff
    \operatorname{card}(b\cap b')=0\ \ \text{\add{et $b,b'$ non associés}}
\end{array}
\]
\note{Le $1$ de la troisième ligne est souligné ; le $3$ de la quatrième
et le $0$ de la dernière sont repassés en gras, le $0$ sur un chiffre
qui pourrait être $4$ ; après chacun, des mots biffés, le dernier
remplacé par l'ajout, au-dessus, « et $b,b'$ non associés ».}

Pour une racine $r$ donnée, il y a
\[
\begin{array}{ll}
  1 & \text{racine égale à}\ r,\ \text{savoir}\ r\\
  1 & \text{racine opposée à}\ r,\ \text{savoir}\ -r\\
  30 & \text{racines orthogonales à}\ r\\
  20 & \struck{\text{\ill{}}}\ \text{faisant un angle de}\ 60^\circ\
    \text{avec}\ r\\
  20 & \text{-----}\ \ 120^\circ\ \ \text{-----}\ \ r
\end{array}
\]
\[
\begin{array}{lll}
  \text{couples de racines} & \text{égales} & 72\\
  & \text{opposées} & 72/2=36\\
  \text{couples de racines orthogonales} && 72\cdot15=1080\\
  \text{couples de racines} & 60^\circ & 72\cdot10=720\\
  & 120^\circ & 72\cdot10=720
\end{array}
\]
Plans des \add{paires de} racines non orthogonales (= plans spéciaux)
$720:6=120$ \struck{\ill{}}

\page{90}
\note{Le haut de la page est barré d'une grande croix et d'un trait
ondulé ; on le donne tel qu'il se lit.}
\struck{$x$ --- $W$} \quad \struck{$x-s_r(x)$}
\[
  s_r(x)=x+(r,x)\,r \qquad s_r(x)-x=(r,x)\,r \qquad
  (r,x)\in\lbrace0,1,-1\rbrace
\]
\[
  (\struck{x-y})\ \ s_rx-x=\qquad \struck{x=wx}\qquad
  wx=x+\rho\quad \rho\ \text{\uncertain{vecteur}}
\]
\[
  s_rwx=s_rx+s_r\rho=x+\underbrace{(x,r)}_{\varepsilon}\,r+s_r\rho
\]
a) $(x,r)=0$ i.e.\ $s_rx=x$ \ill{}

b) $(x,r)=1$ \ill{} \quad $s_rwx-x=s_r\rho+r$
\[
  \struck{r_0-2r_1-r_2+r_6}\qquad
  r_1+2r_2+3r_3+2r_4+r_5+3r_6 \qquad \alpha_1
\]
\note{La croix passe aussi sur la première de ces deux expressions,
qui semble visée ; la seconde n'est touchée que par son bras.}
\[
  3\xi_i-\lambda=3\xi_i-3\eta+\xi=-3\eta+\sum_{\alpha\neq i}\xi_\alpha+4\xi_i
  \qquad 3\uncertain{\eta_i}-\lambda
\]
\[
  \tilde\xi_i=3\xi_i-\lambda=
  \underbrace{3(\xi_i-\xi_1)}_{-3(r_1+r_2+\cdots+r_{i-1})}
  +\underbrace{(\uncertain{3\xi_1-\eta})}_{-r_6+2r_1+r_2}
  +\underbrace{\eta-\lambda}_{-r_0}
\]
\[
  -\tilde\xi_i=3(r_1+\cdots+r_{i-1})\ \struck{\ill{}}\ -r_2+r_6
  +(r_1+2r_2+3r_3)+2r_4+r_5+2r_6
\]
\note{Cette ligne est très surchargée : le début de la parenthèse et
les termes qui suivent $3(r_1+\cdots+r_{i-1})$ sont biffés et repris,
avec un signe « $-$ » souscrit.}
\[
\begin{aligned}
  -\tilde\xi_1&=-r_1+4r_2+3r_3+2r_4+r_5\ \struck{+r_6}\ +3r_6\\
  -\tilde\xi_2&=2r_1+4r_2+3r_3+2r_4+r_5+3r_6\\
  -\tilde\xi_3&=2r_1+4r_2+3r_3+2r_4+r_5+3r_6\\
  -\tilde\xi_4&=2r_1+4r_2+6r_3+2r_4+r_5+3r_6\\
  -\tilde\xi_5&=2r_1+4r_2+6r_3+5r_4+r_5+3r_6\\
  -\tilde\xi_6&=2r_1+\uncertain{9}r_2+6r_3+5r_4+4r_5+3r_6
\end{aligned}
\]
\note{Le coefficient de $r_2$ est, à chaque ligne, un chiffre repassé en
gras sur un autre (on lit $4$, sur $1$ ou $2$ aux deux premières lignes,
et plutôt $9$ à la dernière) ; au $2r_1$ des lignes 2 à 6, le $2$ est
lui aussi repassé.}
\[
  \struck{-\tilde\eta=\ill{}\,r_1+(r_2+r_3)+6r_4+3r_5+9r_6}
\]
\[
  \struck{\tilde\eta=3\eta-3\lambda=-3(\lambda-\eta)=-3r_0}
\]
\[
  -\tilde\eta=+3r_1+6r_2+9r_3+6r_4+3r_5+6r_6
\]

\page{91}
\note{Page entière barrée de deux longs traits ondulés ; elle se lit
d'un bout à l'autre et on la donne sans marquer la rature ligne à
ligne.}
\[
\begin{array}{l}
  a+3b=-2\\
  3a+\uncertain{9}b=0\\
  \struck{2a=}\quad a=-\tfrac{6}{3}\\
  \uncertain{8}a=2\qquad a=\tfrac14\qquad b=
\end{array}
\qquad
\begin{array}{l}
  a+3b+\\
  a+3b+2,\ a+3b+3\\
  a+3b,\ a+3b+1\\
  \quad\parallel\\
  a+3b'-3,\ a+3b'-2
\end{array}
\]
\[
  2\quad -1\qquad a+3b=-1\qquad
  \left.\begin{array}{l}a+3b=-1\\ 3a+b=0\end{array}\right\rbrace\quad
  x=b\bigl(\xi-\tfrac{\uncertain{5}}{3}\bigr)+\xi_i
\]
\[
  x=a\cdot\eta+b\xi+\xi_i\qquad \struck{\ill{}\,6b}\qquad
  x=a\eta+b'\xi-\xi_i\qquad a=2b\qquad \struck{3a+6b}\quad
  \underline{3a+6b=0}
\]

(1) \note{Les numéros (1) à (4) sont cerclés.}
\[
  x=a\eta+b\xi+\xi_i \qquad
  \begin{array}{l}3a+6b=0\\ a+2b=0\\ a=-2b\end{array}
  \qquad
  \begin{array}{ll}
    \alpha)\ a+3b=-1 & \text{i.e.}\ \boxed{b=-1,\ a=2}\quad
      2\eta-\xi+\xi_i=\boxed{\xi'_i}\\
    \beta)\ a+3b=0 & \text{i.e.}\ b=0,\ a=0\quad \boxed{\xi_i}
  \end{array}
\]

(2)
\[
  x=a\eta+b\xi-\xi_i\qquad \lambda_0=\eta-\frac{\xi}{3}
\]
\[
  \text{prenons}\ \ x'=2\lambda_0-x=\struck{\ill{}}\,2\eta-\tfrac23\xi
  -a\eta-b\xi+\xi_i
  =\underbrace{(2-a)}_{a'}\eta
  +\underbrace{\bigl(-\tfrac23-b\bigr)}_{b'}\xi+\xi_i
\]
on est donc ramené au cas \uncertain{précédent} \ill{}

(3)
\[
  x=a\eta+b\xi-\xi_i-\xi_j\qquad
  3a+6b=3\quad a+2b=1\quad a=1-2b\qquad
  a+3b=1\quad \boxed{b=0,\ a=1}\quad \boxed{\xi_{ij}}
\]

(4)
\[
  x=a\eta+b\xi+\xi_i+\xi_j
\]
on se ramène \ill{} \ill{} \ill{} \uncertain{par}
\[
  \delta\sim-(\delta-\lambda_0)+\lambda_0=2\lambda_0-\delta
  =\tfrac23\lambda-\delta=\delta'
\]
\[
\begin{array}{l}
  r_1\longmapsto-r_5\\
  r_2\longmapsto-r_4\\
  r_3\longmapsto-r_3\\
  r_4\longmapsto-r_2\\
  r_5\longmapsto-r_1\\ \hline
  r_6\longmapsto r_{456}=r_6+3(\ill{})+2r_4+r_5
\end{array}
\]
\[
  V_{ijh}=\eta-\xi_i-\xi_j-\xi_h \qquad
  \eta-\xi_3-\xi_5-\xi_{\uncertain{6}}=\eta-\xi_1-\xi_2-\xi_3
  +(\xi_1-\xi_4)+(\xi_2-\xi_5)+\xi_{\ill{}}-\xi_{\ill{}}
\]
\note{Au-dessus de $V_{ijh}$, une fraction $\tfrac23$ ; sous $r_6$, un
renvoi « $r_{12}$ » \uncertain{}. Le bas de la feuille est déchiré et la
fin des deux dernières lignes est perdue.}
\page{92}
\note{Le haut de la page est très raturé ; les deux premières lignes
sont barrées de traits obliques.}
\[
  \xi'_i=\struck{\ill{}}\ r_0+\xi_i \qquad
  \tilde\xi'_i=\tilde r_0+\tilde\xi_i=3r_0+\tilde\xi_i
  \qquad \lbrace i\mapsto\ill{}\rbrace\qquad\lambda
\]
\[
  -\tilde\xi'_i=-\tilde\xi_i-3r_0=-\tilde\xi_i\
  \struck{-3(r_1+r_2+r_3)-6r_4-3r_5-6r_6}\ =-\tilde\xi_i-\tilde\eta
\]
\note{À droite, en colonne : « $\xi'_i-\xi_i=r_0$ », puis
« $3\xi'_i=\struck{\xi_i}\,\ill{}\,3r_0$ », repris en
« $3\xi'_i=3\xi_i+3r_0=3\tilde\xi_i-\tilde\eta$ » et
« $-3\tilde\xi'_i=-\tilde\xi_i+\tilde\eta$ » ; les tildes et le
coefficient $3$ de cette dernière ligne se lisent mal.}
\[
  \xi_{ij}=\eta-\xi_i-\xi_j
\]
\[
  \tilde\xi_{ij}=\tilde\eta-\tilde\xi_i-\tilde\xi_j
  =(-\tilde\xi_i)+(-\tilde\xi_j)-(-\tilde\eta)
\]
\struck{$-\tilde\xi_{ij}=\tilde\xi_i+\tilde\xi_j+\ill{}\,\tilde\eta$}

\ill{}\note{Un mot abrégé, souligné, au-dessus des deux colonnes de
suites ; les deux colonnes sont traversées chacune d'un long trait
oblique.}
\[
\begin{array}{ll}
  1,4,-8,-5,-4,-1 & -8,-5,-4,-1,1,4\\
  2,5,-7,-4,-5,-2 & -7,-5,-4,-2,2,5\\
  3,6,-6,-3,-3,0 & -6,-3,0,3,6\\
  2,5,-4,-1,-2,1 & -4,-2,-1,1,2,5\\
  1,4,-2,1,-1,2 & -2,-1,1,2,4\\
  3,-3,-3 & -3,3
\end{array}
\qquad
\begin{array}{c}6\\6\\5\\6\\5\\2\end{array}
\]
\note{Le $-4$ de la quatrième ligne de gauche est repassé sur un autre
chiffre. La colonne de droite (6, 6, 5, 6, 5, 2), séparée des suites,
est écrite dans la marge.}
\[
  -\tilde\delta=-w(-\tilde\delta)\qquad
  \tilde\delta=\struck{\ill{}}-w(\tilde\delta)=-\widetilde{w(\delta)}
  =-\tilde\delta' \qquad \widetilde{\delta\cdot\delta'}=0
\]
\[
  \ill{}\ \delta+\delta'=2\lambda=\qquad
  \tilde\delta=-w(\tilde\delta)=-\widetilde{w(\delta)}=-\tilde\delta'
  \qquad \tilde\delta\cdot\tilde\delta'=0
\]
Le tableau des coefficients, en bas à gauche :
\[
\begin{array}{c|cccccc}
  & r_1 & r_2 & r_3 & r_4 & r_5 & r_6\\ \hline
  -\tilde\xi_1 & -1 & 1 & 3 & 2 & 1 & 3\\
  -\tilde\xi_2 & 2 & 1 & 3 & 2 & 1 & 3\\
  -\tilde\xi_3 & 2 & 4 & 3 & 2 & 1 & 3\\
  -\tilde\xi_4 & 2 & 4 & 6 & 2 & 1 & 3\\
  -\tilde\xi_5 & 2 & 4 & 6 & 5 & 1 & 3\\
  -\tilde\xi_6 & 2 & 4 & 6 & 5 & 4 & 3\\ \hline
  -\tilde\xi'_1 & 2 & 7 & 12 & 8 & 4 & 9\\
  -\tilde\xi'_2 & 5 & 7 & 12 & 8 & 4 & 9\\
  -\tilde\xi'_3 & 5 & 10 & 12 & 8 & 4 & 9\\
  -\tilde\xi'_4 & 5 & 10 & 15 & 8 & 4 & 9\\
  -\tilde\xi'_5 & 5 & 10 & 15 & 11 & 7 & 9\\
  -\tilde\xi'_6 & 5 & 10 & 15 & 11 & 10 & 9\\ \hline
  & -2 & -4 & -3 & -2 & -2 & 0\\
  & 1 & -1 & 0 & 1 & 1 & 0
\end{array}
\]
\note{Plusieurs chiffres du tableau sont repassés en gras sur d'autres
(le $6$ et le $2$ de la ligne $-\tilde\xi_4$, le $6$ de $-\tilde\xi_5$,
les $10$, $12$, $15$ et $8$ des lignes $-\tilde\xi'_3$ et
$-\tilde\xi'_4$) ; les étiquettes des lignes, écrites serré dans la
marge, se lisent mal et sont données d'après leur ordre.}

En bas à droite :
\[
  -2,-1,1,2,5 \qquad -2,1,4,7,10
\]
\page{95}
\note{Au crayon des archivistes, « 97 » ; voir l'en-tête pour le décalage
de la numérotation à partir d'ici. La page est écrite presque entière à
l'encre bleue pâle, avec des reprises à l'encre plus foncée.}
$R\subset E$ syst.\ de racines ($E$ $\mathbb{R}$-vectoriel engendré par
$R$), produit scalaire invariant\add{ tel que $(r,r)=2$}
\[
  \mathfrak{S}=\bigl\lbrace x\in E^{*}\ \big|\
  (x,r)\ \struck{\tfrac{2}{(r,r)}}\in\lbrace-1,0,1\rbrace\ \
  \forall r\in R\bigr\rbrace
\]
(racines \add{toutes} de \struck{\ill{}} longueur\ldots)\quad
$\to E^{*}=E-\lbrace0\rbrace$
\[
  \mathfrak{S}_0=\bigl\lbrace x\in E^{*}\ \big|\ wx-x\in R\ \
  \forall w\in W\ \text{t.q.}\ wx\neq x\bigr\rbrace
  \qquad W=\text{groupe de Weyl}
\]
\note{La lettre ronde de ces ensembles, une sorte de $\mathfrak{S}$ bouclé,
est rendue $\mathfrak{S}$ partout ; « $W$ » a été écrit sur un « $E$ »
repassé.}

$\mathfrak{S}_0\subset\mathfrak{S}$, car \add{$x\in\mathfrak{S}_0\Rightarrow$}
$\mathfrak{S}_0\in E^{*}$, et $s_rx-x=-2\dfrac{\langle x,r\rangle}{(r,r)}\,r
\in R$, dans le cas où $s_rx\neq x$ i.e.\ $\langle x,r\rangle\neq0$, on a
$-(x,r)\in\dfrac{(r,r)}{2}\lbrace\pm1\rbrace$, donc \struck{\ill{}} dans
tous les cas $(x,r)\in\lbrace-1,0,+1\rbrace$.
\note{« $\mathfrak{S}_0\in E^*$ » est sur la page, pour
« $\mathfrak{S}_0\subset E^*$ ».}

\struck{$\ill{}\ \mathfrak{S}$}

$\operatorname{card}\mathfrak{S}\leq3^r-1$ \quad donc \ill{} $\mathfrak{S}$
\ill{} a fortiori $\mathfrak{S}_0$ fini.
$\mathfrak{S}$, $\mathfrak{S}_0$ stables par $W$ \ill{} $=$ \ill{}
$\hat W=\operatorname{Aut}(E,R)$, et $W$ a un nb fini d'orbites dans
$\mathfrak{S}$ et \ill{} fortiori dans $\mathfrak{S}_0$.

Cas $(A_r)$\note{L'indice de $A$ est repassé en gras et se lit mal.}
\[
  E=\mathbb{R}^{r+1},\qquad
  E=\Bigl\lbrace(x_1,\ldots,x_{r+1})\ \Big|\ \sum_{1}^{r+1}x_i=0\Bigr\rbrace
\]
\[
  R=\lbrace\alpha_{ij}=\xi_i-\xi_j\mid i\neq j,\ 1\leq i,j\leq r+1\rbrace
  \qquad \operatorname{card}R=r(r+1)
\]
produit scalaire \uncertain{induit} par
$(\xi,\eta)=\sum_1^{r+1}\xi_i\eta_i$ sur $\mathbb{R}^{r+1}$
\struck{\ill{}}
\[
  \xi_1\ \ldots\ \xi_r\qquad \sum\xi_i=0
\]
\[
  \underbrace{\xi_1=\cdots=\xi_{r_1}}_{\alpha_1}<
  \underbrace{\xi_{r_1+1}=\cdots=\xi_{r_1+r_2}}_{\alpha_2}\ \cdots\
  \underbrace{\xi_{r_1+r_2+\cdots+r_{s-1}+1}=\cdots=
    \xi_{\underbrace{r_1+r_2+\cdots+r_s}_{r+1}}}_{\alpha_s}
\]
\[
  \begin{cases}
    r_1\alpha_1+r_2\alpha_2+\cdots+r_s\alpha_s=0\\
    r_1+r_2+\cdots+r_s=r+1
  \end{cases}
  \qquad
  \begin{array}{l}
    \alpha_1<\alpha_2<\cdots<\alpha_s\\
    r_1,r_2,\ldots,r_s\geq1
  \end{array}
\]
\[
  \alpha_i-\alpha_j\in\lbrace\pm1\rbrace\quad\forall i,j\ \neq
\]
\note{Le « $\neq$ » final, isolé, précise sans doute « $i\neq j$ ».}

Soit \struck{\ill{}} $s=1$, \struck{\ill{}} $r_1=r+1$,
$r_1\alpha_1=(r+1)\alpha_1=0$, $\alpha_1=0$, absurde ($\xi=0$).

Soit $s=2$, $\alpha_2=\alpha_1+1$
\[
  \begin{cases}
    r_1\alpha_1+r_2\alpha_2=0\\
    r_1+r_2=r+1\\
    \alpha_2=\alpha_1+1
  \end{cases}
  \Longrightarrow
  \begin{array}{l}
    (r_1+r_2)\alpha_1+r_2=0\\
    \text{i.e.}\ (r+1)\alpha_1+r_2=0\\
    \text{i.e.}\ \alpha_1=\dfrac{-r_2}{r+1}
  \end{array}
\]
\[
  \struck{\text{Il y a}\ r+}\quad
  \alpha_1=\frac{-r_2}{r+1},\quad\alpha_2=\frac{r_1}{r+1}\qquad
  r_1+r_2=r+1
\]
il y a $r+1$ classes de tels éléments, suivant la valeur de
$1\leq r_1\leq r+1$\note{Lecture douteuse de la fin de la phrase ; la
borne supérieure est sur la page « $r+1$ ».}
\page{96}
\note{Au crayon, « 98 ». Feuille de listing tenue en largeur, en deux
colonnes : à gauche la fin du cas $A_r$ et le cas $B_r$ ; à droite les
cas $C_r$ et $D_r$. La suite de la page 95.}
\[
  \underbrace{-\frac{r_2}{r+1},\ldots,-\frac{r_2}{r+1}}_{r_1},\
  \underbrace{\frac{r_1}{r+1},\ldots,\frac{r_1}{r+1}}_{r_2}
  \qquad
  -\frac{r_1}{r+1}\ \ -\frac{r_2}{r+1}\qquad\pm1
\]
\[
  [1,r+1]=\underbrace{[1,r_1]}_{I_1}\cup
  \underbrace{[r_1+1,r+1]}_{I_2}
  \qquad \sigma I_1,\ \sigma I_2
\]
$\xi$ \quad $\alpha_1$ sur $I_1$, $\alpha_2$ sur $I_2$

$\sigma\xi=\xi\circ\sigma^{-1}$ \quad $\alpha_1$ sur $\sigma I_1$, $\alpha_2$
sur $\sigma I_2$
\[
  \xi-\sigma\xi\qquad
  \begin{array}{ll}
    0 & \text{sur}\ (I_1\cap\sigma I_1)\cup(I_2\cap\sigma I_2)\\
    \struck{-1} & \text{sur}\ I_1\cap\sigma I_2\\
    & \text{sur}\ I_2\cap\sigma I_1
  \end{array}
\]
On veut, $\forall\sigma$ tel que $\sigma I_1\neq I_1$ \struck{\ill{}}
i.e.\ $\sigma I_1\cap I_2\neq\emptyset$ \add{i.e.\
$\sigma I_2\cap I_1\neq\emptyset$} et $\sigma I_2\neq I_2$
\struck{$I_1\cap\sigma I_2$}
\[
\begin{array}{ll}
  \operatorname{card}(I_1\cap\sigma I_2)=1 &
    \to\ \text{c'est possible (\uncertain{ssi}}\
    \operatorname{card}I_1\leq1\ \text{ou}\ \operatorname{card}I_2\leq1\
    \forall\sigma\ldots)\\
  \operatorname{card}(I_2\cap\sigma I_1)=1 & \text{--- \,--- }
\end{array}
\]
\struck{Donc} OK ssi $r_1=1$ ou $r_2=1$
\[
  \underbrace{-\frac{1}{r+1},\ \frac{r}{r+1},\ldots,\frac{r}{r+1}}_{r}
  \qquad
  \underbrace{-\frac{r}{r+1},\ \struck{\ill{}},\ -\frac{r}{r+1}}_{r},\
  \frac{1}{r+1}
\]
\note{Les accolades sont sur la page ; le « $r$ » de la première ne
couvre que les termes $\frac{r}{r+1}$.}
si $r\geq2$, 2 orbites sous $W$, qui sont \struck{\ill{}}
\add{échangées} l'une de l'autre par \struck{\ill{}} l'élément
\struck{\ill{}} \ill{} de $E=\hat W/W$ i.e.\ par la symétrie !


\textbf{($B_r$)}\note{Les étiquettes $B_r$, $C_r$, $D_r$ sont
cerclées.}\quad $V=E=\mathbb{R}^r$
\[
  R=\lbrace\pm\varepsilon_i\ (1\leq i\leq r),\ \pm\varepsilon_i\pm
  \varepsilon_j\ (1\leq i<j\leq r)\rbrace
  \qquad 2r+4\,\frac{r(r-1)}{2}=2r+2r(r-1)=2r^2
\]
\drawing{À droite de $V=E$, un croquis de matrice d'ordre $2r+1$ : des
blocs $a_{ij}$ et $-a_{ij}$ de part et d'autre de la diagonale, un $0$
en bas.}
Produit scalaire habituel, longueurs $1$ et $2$
\[
  \xi=(\xi_1,\xi_2,\ldots,\xi_r)\qquad
  \struck{\xi_i}\ \ 0\leq\xi_1\leq\xi_2\leq\cdots\leq\xi_r
\]
\[
  \underbrace{\xi_1=\cdots=\xi_{r_1}}_{\alpha_1}\
  \underbrace{\xi_{r_1+1}=\cdots=\xi_{r_1+r_2}}_{\alpha_2}\ \cdots\
  \underbrace{\xi_{r_1+\cdots+r_{s-1}+1}=\cdots=
    \xi_{r_1+\cdots+r_s=r}}_{\alpha_s}
\]
\[
  (\xi,\varepsilon_i)\in\struck{2}\lbrace-1,0,1\rbrace\ \ \text{d.c.}\ \
  \boxed{\xi_i\ \text{égale}\ 0\ \text{ou}\ 2}
\]
\note{Dans l'encadré, le mot « égale » est incertain et le « 2 » est
peut-être la fin d'un « $\pm1$ » ; le $2$ de la ligne est biffé. La
ligne suivante et la conclusion $\xi=(2,2,\ldots,2)$ ne sont pas
cohérentes avec « $\xi_i=0$ ou $\pm1$ » : on transcrit tel quel.}
\[
  (\xi,\varepsilon_i\pm\varepsilon_j)=\xi_i\pm\xi_j\in\lbrace1,0,-1\rbrace
  \qquad \struck{\ill{}}\
  \begin{array}{l}\alpha_i+\alpha_j=1\\ \alpha_i-\alpha_j=-1\end{array}
  \qquad \struck{\ill{}}\ \ \text{\uncertain{si}}\ 1\leq i<j\leq s
\]
\struck{Une seule} donc impossible, donc $s=1$
\[
  \xi=(2,2,\ldots,2)\qquad \mathfrak{S}=\emptyset
\]

\textbf{($C_r$)}\quad $E=\mathbb{R}^r$
\[
  R=\lbrace\pm2\varepsilon_i\ (1\leq i\leq r),\ \pm\varepsilon_i\pm
  \varepsilon_j\ (1\leq i<j\leq r)\rbrace
\]
\[
  \struck{\xi_i\ \ldots}\ \ (\xi_i)_{1\leq i\leq r}\qquad
  \xi_i\in\tfrac24\lbrace-1,0,+1\rbrace=\lbrace\ \rbrace
\]
\[
  0\leq\xi_1\leq\cdots\leq\xi_r\quad\text{i.e.}\quad
  \xi_1\in\lbrace0,\tfrac12\rbrace\qquad
  \pm\xi_i\pm\xi_j\in\lbrace-1,0,+1\rbrace
\]
\[
  \struck{\ill{}}\ \bigl(\tfrac12,\tfrac12,\ldots,\tfrac12\bigr)\qquad
  \struck{r\geq3}\ \ r\geq3
\]
\[
  \mathfrak{S}=\struck{\ill{}}\ \struck{\varepsilon_1,\varepsilon_2}\
  \ill{}\,\lbrace\pm\tfrac12\rbrace^{r}\qquad
  \mathfrak{S}_0=\emptyset
\]

\textbf{($D_r$)}\quad $E=\mathbb{R}^r$,\quad
$R=\lbrace\pm\xi_i\pm\xi_j\ (1\leq i<j\leq r)\rbrace$
\[
  \struck{\ill{}}\ \xi_1\ \ldots\ \xi_r\qquad
  \xi_1\leq\xi_2\leq\cdots\leq\xi_r\qquad
  \pm\xi_i\pm\xi_j\in\lbrace-1,0,1\rbrace
\]
\note{Sous $\xi_2$, un renvoi « $\geq0$ » \uncertain{}. Au-dessous, un
croquis de matrice diagonale $\operatorname{diag}(\lambda_1,
\lambda_1^{-1},\ldots,\lambda_r,\lambda_r^{-1})$, entouré d'une courbe
qui la relie à la colonne de droite.}
au plus deux valeurs distinctes des $\xi_i$, différant de $1$, soit
$\alpha$ et $\alpha+1$, mais \struck{$\alpha+$} $2\alpha+1\in\lbrace-1,0,1
\rbrace$ i.e.\ \ill{} \ill{} les valeurs égales, égales : $\pm\frac12$
\[
  2\alpha+1=0\quad\alpha=-\tfrac12\quad\alpha+1=\tfrac12\qquad
  2\alpha+1=1,\ \alpha=0\qquad 2\alpha+1=-1,\ \alpha=-1
\]
\[
  \left\lbrace
  \begin{array}{l}
    -\frac12,\ \underbrace{\frac12,\frac12,\ldots,\frac12}_{r-1}\\[4pt]
    (\underbrace{0,\ldots,0}_{r-1},1)\ \ \struck{\ill{}}\
      \boxed{r_2=1}\ \longrightarrow r\\[4pt]
    (-1,\underbrace{0,\ldots,0}_{r-1})\ \ \struck{\ill{}}\
      \text{\uncertain{conjugués}}\ (r_2=1)\\[4pt]
    (\frac12,\frac12,\ldots,\frac12)
  \end{array}\right.
\]
\note{Une flèche descend de la deuxième ligne à la troisième ; à côté,
biffés, des mots illisibles.}
\[
  \bigl(\tfrac12\ \struck{\ill{}}\bigr)\qquad \struck{\ill{}}\qquad
  \prod(\ill{})=1\qquad \prod(\ill{})=-1\qquad
  \struck{\ill{}\,\lbrace\pm\varepsilon_i\rbrace}
\]
il y a \ill{}\quad $\mathfrak{S}$ a trois orbites,\quad
$\mathfrak{S}_0=\emptyset$\ \struck{\ill{}}
\page{97}
\note{Au crayon, « 99 ». Reprise des définitions de la page 95.}
\[
  \mathfrak{S}'_0=\bigl\lbrace x\in E\ \big|\ w\in W,\
  w(\struck{\ill{}}\,\mathbb{R}x)\neq\mathbb{R}x\Longrightarrow
  wx-x\in R\bigr\rbrace
\]
\[
  x\in\mathfrak{S}'_0,\ r\in R\qquad
  s_rx-x=2\,\frac{(x,r)}{(r,r)}\,r
\]
\[
  s_r(\struck{\ill{}}\,\mathbb{R}x)=(\mathbb{R}x)\iff
  s_rx=x\ \text{ou}\ s_rx=-x\iff(x,r)=0\ \text{ou}\ x\in\mathbb{R}r
\]
\[
  (x,r)\neq0\ \text{et}\ x\notin\mathbb{R}r\Longrightarrow
  2\underbrace{\frac{(x,r)}{(r,r)}}_{\uncertain{n_{x,r}}}=\pm1
\]
\struck{$(x,r)$} donc
\[
  \mathfrak{S}'_0\subset\mathfrak{S}'=\Bigl\lbrace x\in E\ \Big|\
  \forall r\in R\cap\complement\,\mathbb{R}x,\ \text{on a}\
  (x,r)\in\frac{(r,r)}{2}\lbrace-1,0,+1\rbrace\Bigr\rbrace
  \qquad
  \begin{array}{ccc}
    \mathfrak{S}'_0&\subset&\mathfrak{S}'\\
    \cup&&\cup\\
    \mathfrak{S}_0&\subset&\mathfrak{S}
  \end{array}
\]
\note{Le $\complement$ est une grande lettre repassée en gras.}
\[
  \mathfrak{S}'=\mathfrak{S}\cup\Bigl\lbrace\lambda s\ \Big|\ s\in R,\
  \lambda\in\mathbb{R}^{*+},\ (\lambda s,r)\in\frac{(r,r)}{2}
  \lbrace-1,0,1\rbrace\ \ \forall r\in R-\lbrace s,-s\rbrace\Bigr\rbrace
\]
\[
  \text{i.e.}\ \ \forall r\in R-\lbrace s,-s\rbrace\ \struck{\ill{}}\
  (R\cap s^{\perp})\qquad
  \struck{\text{i.e.}\ \lambda\in\frac{(r,r)}{2(s,r)}}\qquad
  \lambda\in\frac{(r,r)}{2(s,r)}\lbrace-1,0,1\rbrace
\]
\[
  \text{i.e.}\ \ \lambda=\pm\frac{(r,r)}{2(s,r)}\qquad
  \text{i.e.}\ \ \frac1\lambda=\pm\frac{2(s,r)}{(r,r)}=\pm n(s,r)
\]
\note{La ligne « i.e.\ $\forall r\in R-\lbrace s,-s\rbrace$ » précise
l'ensemble placé sous l'accolade ; le signe entre $\lbrace s,-s\rbrace$
et $(R\cap s^\perp)$ est surchargé, peut-être un « $-$ ».}
Un tel $\lambda$ existe (il est unique au signe près), \uncertain{pour}
$s$ donné, ssi les quantités $\Bigl|\dfrac{(s,r)}{(r,r)}\Bigr|$, pour
$r$ racine non orthogonale ni \uncertain{proportionnelle} à $s$, sont
égales, i.e.\ ssi il n'y a \add{qu'un seul} \struck{\ill{}} angle de
droites possible (\uncertain{savoir} $\frac\pi3$) \add{d'où
$\uncertain{n}(s,r)=\pm1$, d'où $\lambda=\pm1$} pour des racines
distinctes non orth. C'est le cas ssi toutes les racines sont de même
longueur, i.e.\ les types $A_r$, $D_r$, $E_i$ ($6\leq i\leq8$).
\note{Les deux ajouts sont écrits au-dessus de la ligne, reliés par un
trait à « il y a » ; le second est en partie masqué.}

\page{98}
\note{Au crayon, « 100 ».}
\[
  \mathfrak{S}'_0=
  \begin{cases}
    \mathfrak{S}_0 & \text{si toutes les racines ne sont pas de même
      longueur}\\
    \mathfrak{S}_0\cup R & \text{si toutes les racines de même long.}
  \end{cases}
\]
\note{Au-dessous, un bloc de trois lignes biffé de hachures et d'un
cadre : « $\mathfrak{S}_0$ \ill{} si toutes les \ill{} »,
« $\mathfrak{S}'_0\ \ill{}\ \mathfrak{S}_0$ si toutes les racines sont
de même longueur », « $\mathfrak{S}'_0\cap\mathfrak{S}$ ».}

\emph{Cas de $A_r$} ($r\geq2$)\quad $\mathfrak{S}'=\mathfrak{S}\cup R$
\[
  \text{\uncertain{rel.\ donnés}}\quad
  \struck{\underbrace{-\tfrac{r_2}{r+1},\ \ldots,\ -\tfrac{r_2}{r+1}}}
  \quad
  \left\lbrace
  \begin{array}{l}
    -\frac{r_2}{r+1}(e_1+\cdots+e_{r_1})+\frac{r_1}{r+1}
      (e_{r_1+1}+\cdots+e_{r_1+r_2})\\
    \quad 1\leq r_2\leq r-1\\
    e_1-e_2
  \end{array}\right.
\]
\[
  \mathfrak{S}'_0=
  \begin{cases}
    \mathfrak{S}_0 & \text{si}\ r\neq3\ \text{i.e.}\ r+1\neq4\\
    \struck{\ill{}}\ \mathfrak{S}_0\cup W.\bigl(-\tfrac12,-\tfrac12,
      \tfrac12,\tfrac12\bigr)
  \end{cases}
\]
\note{Au-dessus de « $r\neq3$ », un segment coupé en deux, marqué $r_1$,
$r_2$.}

\page{100}
\note{Au crayon, « 102 ». Une page de théorie des groupes, sans lien
apparent avec ce qui précède.}
\[
  G\qquad N_1\ N_2\ N_3
\]
\drawing{Un petit diagramme : $G$ au centre, avec une flèche vers le
haut vers $G/N_1$, une vers la droite vers $G/N_2$, et un trait vers le
bas à gauche vers $G/N_3$.}
\[
  N_1\cap N_2=\lbrace e\rbrace\qquad
  \forall x_1\in G,\ x_2\in G\quad \exists x\in G,\quad
  x\in x_1N_1\cap x_2N_2\quad
  \text{i.e.}\ \begin{cases}x_1^{-1}x\in N_1\\ x_2^{-1}x\in N_2\end{cases}
\]
\[
  N_1\times N_2\to G\qquad(x_1,x_2)\mapsto x_1x_2\qquad
  \struck{N_1N}\ N_1.N_2=G\iff N_2N_1=G
\]
\[
  x_1x_2=y_1y_2\qquad(y_1^{-1}x_1)=(y_2x_2^{-1})\in N_1\cap N_2=\lbrace
  e\rbrace\qquad x_1=y_1,\ x_2=y_2
\]
\[
  \struck{x_1\in G}\qquad a_1b_2=x_1\quad d_2c_1=x_2\qquad a_1\ \ a
\]
\[
  N_1\times N_2\simeq G\qquad
  \begin{array}{l}N_1\simeq G/N_2\\ N_2\simeq G/N_1\end{array}
\]
\note{De « $N_1\times N_2\simeq G$ » partent deux flèches croisées vers
$N_1\simeq G/N_2$ et $N_2\simeq G/N_1$, et une longue flèche vers un
$x$ isolé à droite.}
\[
  \struck{(x_1,x_2)\mapsto x_1x_2}\qquad
  \begin{array}{l}
    x=a_1a_2\ \longrightarrow\ a_1\in N_1\simeq G/N_2\qquad a_1=x_1\\
    \phantom{x}=b_2b_1\ \longrightarrow\ \struck{\ill{}}\ \
      N_2\simeq G/N_1
  \end{array}
\]
\note{Entre ces deux lignes, biffé, « $x_1x_2\to x_1$ ».}
\[
  x=x_1a_2\qquad x=x_2b_1\qquad x_1a_2=x_2b_1\qquad
  x_2^{-1}x_1=b_1a_2^{-1}=b_1a'_2
\]
\[
  \begin{cases}
    N_2\cap N_3=N_3\cap N_1=N_1\cap N_2=\lbrace e\rbrace\\
    N_2N_3=N_3.N_1=N_1.N_2=G
  \end{cases}
  \qquad
  G\simeq N_2\rtimes N_1\simeq N_3\rtimes N_1
\]
\note{Les deux produits semi-directs sont écrits comme un « $\cdot$ »
surmonté d'un petit signe ; la lecture $\rtimes$ est incertaine.}
À droite, séparé par un trait : $G$ \ill{}, $N_1$ invariant,
$G/N_1\simeq V$ \uncertain{} ; $N_2$ \ \ $N_3$.
\end{document}
